% =====================================================================
\section{Related Work}\label{sec:related}
Coordination of autonomous agents has been studied for more than three decades, and almost every mechanism in this paper has an ancestor. We searched the main venues of the communities involved: software engineering (ICSE, FSE, ASE, ISSTA, TSE, TOSEM), machine learning and NLP (NeurIPS, ICML, ICLR, ACL), multi-agent systems (AAMAS, JAAMAS, EMAS, EUMAS), and databases, programming languages and systems (SIGMOD, VLDB, PLDI, NSDI). Because the area moves quickly, we also include preprints that appeared up to November 2026 and mark them as concurrent work. Table~\ref{tab:positioning} summarizes the comparison along the five properties our controller combines. The section closes by stating which of these properties we borrow and which we add.

\begin{table*}[!t]
\centering
\caption{Positioning. \emph{C1}: a contribution is admitted only if it was computed from inputs that are still current, judged by a trusted component; \emph{C2}: invalidation is exact and content-based, with early cutoff; \emph{C3}: retries are bounded per dependency episode, with a terminal-outcome guarantee; \emph{C4}: publication is gated by a global consistency predicate; \emph{C5}: the guarantees are proved for a task-independent coordinator and machine-checked. $\bullet$: provided; $\circ$: partially or in a different form (see text); --: not addressed. All families except build systems and transactions admit autonomous, untrusted participants.}\label{tab:positioning}
\footnotesize\setlength{\tabcolsep}{4pt}
\begin{tabular*}{\textwidth}{@{\extracolsep{\fill}}p{0.30\textwidth}p{0.28\textwidth}ccccc@{}}\toprule
Family & Representative work & C1 & C2 & C3 & C4 & C5\\\midrule
Tuple centres, agents \& artifacts & \cite{omicini2001tuple,omicini2008artifacts,ricci2011environment} & $\circ$ & -- & -- & $\circ$ & $\circ$\\
Commitment and information protocols & \cite{singh1998acl,yolum2002commitments,singh2011bspl,singh2020clouseau} & $\circ$ & -- & -- & -- & $\circ$\\
Electronic institutions, normative MAS & \cite{esteva2001institutions,arcos2005institutions,boella2006normative} & $\circ$ & -- & -- & $\circ$ & --\\
MAS model checking & \cite{lomuscio2017mcmas,wooldridge2002mable,dennis2012ajpf} & -- & -- & -- & -- & $\circ$\\
LLM multi-agent frameworks and orchestrators & \cite{metagpt,chatdev,autogen,langgraph,zhuge2024gptswarm,zhang2025aflow} & -- & -- & $\circ$ & -- & --\\
Runtime enforcement for LLM agents & \cite{wang2026agentspec,chen2025shieldagent,kamath2025agentc,debenedetti2025camel} & $\circ$ & -- & -- & -- & --\\
Transactions for LLM agents (concurrent) & \cite{chang2025sagallm,lyu2026coagent,li2026memtx,chen2026planfence} & $\circ$ & $\circ$ & $\circ$ & -- & --\\
Verified LLM-agent protocols and workflows & \cite{xia2026tracefix,wang2026lean4agent} & -- & -- & -- & $\circ$ & $\circ$\\
Build systems, incremental computation & \cite{mitchell2012shake,mokhov,erdweg2015pluto,hammer2014adapton} & $\circ$ & $\bullet$ & -- & -- & $\circ$\\
Optimistic concurrency control & \cite{occ,berenson1995critique} & $\circ$ & -- & -- & -- & --\\
\midrule
This paper & & $\bullet$ & $\bullet$ & $\bullet$ & $\bullet$ & $\bullet$\\
\bottomrule\end{tabular*}
\end{table*}

% ---------------------------------------------------------------------
\subsection{Coordination in Multi-Agent Systems}
\emph{Coordination models and artifacts.} Malone and Crowston define coordination as managing dependencies between activities~\cite{malone1994coordination}. Their dependency types (prerequisite, shared resource, simultaneity) correspond to our ranked prerequisites, owned slots and atomic transactions. Jennings identifies commitments and conventions as the foundation of coordination~\cite{jennings1993commitments}, and GPGP coordinates task-structured agents by detecting relationships between their tasks~\cite{decker1995gpgp}. Coordination languages separate computation from coordination~\cite{gelernter1992coordination}. In Linda, agents communicate only through a shared tuple space~\cite{gelernter1985linda}. Tuple centres such as ReSpecT and TuCSoN make the medium's reactions programmable, so coordination laws are enforced by the medium, not by the agents~\cite{omicini2001tuple,omicini2001coordination}. The A\&A meta-model makes the environment a set of first-class artifacts with usage interfaces and their own guards~\cite{omicini2008artifacts,weyns2007environment}, realized in CArtAgO and JaCaMo~\cite{ricci2011environment,boissier2013jacamo}. We inherit three ideas from this lineage: \emph{objective coordination}, in which the laws live in a trusted medium; artifacts as the only interface between agents; and guarded operations. Our controller is an A\&A artifact whose single agent-facing operation is ``reply to request $u$''. Generative proposals need a different notion of admissibility. A tuple-centre reaction or an artifact guard decides from the \emph{current} state whether an operation may take effect (hence $\circ$ for C1). Our guard must also decide whether the \emph{content} of the operation was computed from inputs that are still current. That requires recording the observation at dispatch, exact dependency keys and per-episode identities. Tuple centres also have no notion of minimal, content-based invalidation with early cutoff, and this is what makes coordination affordable when every operation costs an LLM call. The EMAS community has begun to relate these abstractions to LLM-based agents~\cite{ciortea2025eumas}; our controller is one concrete answer to how an artifact should govern agents whose outputs are generated rather than computed.

\emph{Protocols, commitments and institutions.} Singh argues that agent communication should have a public, social semantics that can be checked without access to agents' mental states~\cite{singh1998acl}. Commitment protocols operationalize this~\cite{yolum2002commitments}, and Clouseau generates protocols from commitment specifications~\cite{singh2020clouseau}. BSPL specifies a protocol by the information each message carries, identified by keys, and derives causality and the absence of conflicting bindings from those keys~\cite{singh2011bspl,chopra2020protocols}. Our request identities and dependency keys play a similar role. A reply is meaningful only for the binding of keys under which it was requested, and two values for the same slot and episode cannot both be committed. Electronic institutions use governors that let an agent's message take effect only if the institution's current state permits it~\cite{esteva2001institutions,arcos2005institutions}, and normative MAS treat obligations and permissions as first-class runtime state~\cite{boella2006normative}. Our controller is such a governor. The outstanding obligations $O(S)$ play the role of active obligations, and the ownership and authority checks are permissions. The difference lies in what is governed. Institutions govern the \emph{admissibility of actions in a protocol}. We govern the \emph{currentness and validity of generated content} and the \emph{completion} of a set of interdependent artifacts. The two are complementary: a protocol language such as BSPL could specify the interaction among agents, and our controller would decide which of the resulting artifacts are usable and publishable. Interoperability protocols for LLM agents (MCP, A2A, ACP, ANP) standardize transport, discovery and tool access~\cite{mcp,ehtesham2025protocols}, but none of them specifies when a returned artifact may be used.

\emph{Verification of multi-agent systems.} Dedicated MAS model checkers verify temporal, epistemic and strategic properties of agent programs or interpreted systems. MCMAS checks CTLK and ATL specifications of interpreted systems~\cite{lomuscio2017mcmas}, MABLE and AJPF model-check BDI-style agent programs~\cite{wooldridge2002mable,bordini2006verifying,dennis2012ajpf}, and Dastani et al.\ survey the field~\cite{dastani2010verification}. These tools fit when the agents' programs or knowledge are the object of verification. In our setting the agents are opaque, nondeterministic environment components. The object of verification is a single deterministic coordinator with data-rich state: unbounded identifiers, keys, sets of requests and histories. Its properties are state invariants and a refinement between an asynchronous and an atomic machine. TLA$^+$ and its tools were designed for exactly such properties and are routinely used for them in industry~\cite{newcombe2015aws,howard2025smartcasual}. An epistemic formulation (``the controller knows that the value is current'') would be natural in MCMAS and is an interesting complement, but it would require a finite abstraction of keys and identities.

% ---------------------------------------------------------------------
\subsection{LLM-Based Agents and Their Coordination}
\emph{Multi-agent software engineering.} Role-based frameworks encode development processes as conversations among LLM agents. MetaGPT~\cite{metagpt} encodes standard operating procedures as role prompts and structured hand-offs. ChatDev~\cite{chatdev} organizes agents into a chat chain, and AutoGen~\cite{autogen}, CAMEL~\cite{li2023camel} and AgentVerse~\cite{chen2024agentverse} provide programmable multi-agent conversation. For code generation, Self-collaboration assigns analyst, coder and tester roles~\cite{dong2024selfcollab}. AgentCoder separates a programmer from an independent test designer and executor~\cite{huang2023agentcoder}, and MapCoder chains retrieval, planning, coding and debugging agents~\cite{islam2024mapcoder}. Repository-level agents such as SWE-agent~\cite{sweagent}, OpenHands~\cite{wang2025openhands} and AutoCodeRover~\cite{zhang2024autocoderover} are evaluated on SWE-bench~\cite{swebench}. Agentless shows that a fixed localize-repair-validate pipeline is competitive with autonomous agents~\cite{xia2025agentless}, which supports our premise that a deterministic workflow with LLM leaves is a useful design point. Surveys map the landscape~\cite{hou2024llm4se,he2025mas,liu2024agentsurvey,wang2024agents}. All these systems aim to improve \emph{what} the agents produce. They pass artifacts by message or through a shared pool and implicitly assume that what an agent read is still current when its output is used. Our work makes that assumption explicit and enforces it, and could host any of these agents as proposers.

\emph{Workflow orchestration and optimization.} Stateful graph orchestrators such as LangGraph~\cite{langgraph} offer typed shared state with reducer functions, checkpointing and conditional edges. DSPy compiles declarative pipelines of LLM calls~\cite{khattab2024dspy}, and GPTSwarm and AFlow represent agent systems as graphs and optimize their prompts, edges or code~\cite{zhuge2024gptswarm,zhang2025aflow}. These systems provide the \emph{mechanisms} our controller needs (typed state, deterministic reducers, checkpoints, graph structure), but they leave the \emph{policy} to the application. A reducer applies whatever update arrives. Nothing records the observation an update was computed from, rejects an update whose inputs have changed, invalidates exactly the affected state after an edit, or bounds retries per dependency episode (hence $\circ$ for C3: iteration limits exist, but per run). The rules of Section~\ref{sec:implications} can be implemented as such a policy, for example as a LangGraph reducer that commits a node's output only if its recorded key is current. Workflow optimizers are complementary: they choose the graph, and our controller keeps a chosen graph consistent under change.

\emph{Feedback and self-repair.} Iterative refinement with self-feedback~\cite{madaan2023selfrefine}, verbal reinforcement~\cite{shinn2023reflexion}, self-debugging with execution feedback~\cite{chen2024selfdebug}, selection by generated tests~\cite{chen2023codet} and re-prompting with failing tests~\cite{schafer2024testpilot} all show that validator diagnostics improve the next attempt. Our bounded retries with feedback are an instance of this idea. What we add is where the bound lives. A per-run iteration cap does not survive concurrent edits, because a retry may be answering a question that has since changed. A budget per dependency episode does survive them, and it is what makes the terminal-outcome and single-episode theorems hold.

\emph{Failures of multi-agent LLM systems.} Cemri et al.\ annotate more than 1,600 traces from seven frameworks and report inter-agent misalignment and missing task verification as major failure categories~\cite{cemri2025mast}. CooperBench finds that two cooperating coding agents succeed markedly less often than one agent doing both tasks~\cite{khatua2026cooperbench}. Xia et al.\ construct parallel patches that pass alone and fail together because one agent changed an interface the other still relied on~\cite{xia2026stale}. These studies characterize the problem empirically. Our observation keys are a mechanism against the stale-dependency failures they describe, and our global predicate against the jointly inconsistent ones.

% ---------------------------------------------------------------------
\subsection{Guarding, Transacting and Verifying LLM Agents}
\emph{Runtime enforcement for LLM agents.} AgentSpec lets users write trigger-predicate-enforcement rules over agent actions and enforces them at run time~\cite{wang2026agentspec}. ShieldAgent extracts verifiable policies from regulations and checks action trajectories against them~\cite{chen2025shieldagent}. Agent-C enforces temporal constraints on tool calls with an SMT solver during decoding~\cite{kamath2025agentc}. CaMeL separates a privileged planner from a quarantined model and tracks data provenance through a custom interpreter to defeat prompt injection~\cite{debenedetti2025camel}, and AgentDojo evaluates such defenses~\cite{debenedetti2024agentdojo}. These systems guard \emph{actions}: each proposed step is admitted or blocked by a predicate over the current trace. Our guard admits \emph{artifacts} by a predicate over the observation they were computed from, and it records how admitted artifacts depend on each other. Provenance tracking in CaMeL is the closest idea (hence $\circ$), but it serves confidentiality and integrity, not currentness. The approaches compose: an action guard could sit between a proposer and its tools, and our controller would sit between proposers and the shared workspace. Classical runtime verification and enforcement study the same distinction between observing and constraining executions~\cite{schneider,leucker2009rv,pinisetty}; shields correct reactive outputs to satisfy safety specifications~\cite{shield}, whereas our controller defers or rejects proposals instead of correcting them.

\emph{Transactions and concurrency control for LLM agents (concurrent work).} Several recent proposals bring database ideas to agents. Yang et al.\ argue that multi-agent systems should prioritize concurrency control, because long inference windows amplify stale reads and lost updates~\cite{yang2026concurrency}. SagaLLM applies the saga pattern~\cite{garcia1987sagas} to multi-agent planning, with compensating actions and external validation of inter-agent dependencies~\cite{chang2025sagallm}. CoAgent fixes a serialization order at launch, serves order-filtered reads and lets each agent's LLM judge whether a conflicting write invalidates its plan~\cite{lyu2026coagent}. Mraz et al.\ call for orchestrators that track each agent's \emph{state footprint}~\cite{mraz2026footprints}, a notion close to our read set and observation. MemTX stages agent-memory writes through validity levels and repairs derived memory when a belief is retracted~\cite{li2026memtx}, which resembles our separation of accepted from usable values and our invalidation. PlanFence is the closest of these mechanisms. Before a protected action, it follows application-supplied dependency links to the action's inputs, refreshes those that changed and allows one replan before blocking~\cite{chen2026planfence}. Our design differs from this group in four ways. (i) The staleness judgment stays in the trusted controller and is computed from exact keys that include the instruction, validator and selected context, not only the data read; delegating it to the producing model, as in CoAgent, would trust the component whose outputs are in question. (ii) Invalidation is content-based, so equal regenerations stop propagation (early cutoff), which the version- or log-based schemes do not provide. (iii) Retries are bounded per dependency episode, with a proved terminal outcome; the one-replan rule of PlanFence is a bound of this kind for a single action. (iv) Nothing becomes externally visible before the publication barrier, so no compensation is needed, and the guarantees are proved and partly mechanized for a task-independent kernel. Several of these works report empirical evaluations; none, to our knowledge, states and proves safety, termination and cost theorems for its protocol.

\emph{Verified protocols and workflows for LLM agents.} TraceFix generates PlusCal coordination protocols for LLM agents and repairs them with TLC counterexamples until verification passes~\cite{xia2026tracefix}. Lean4Agent models agent workflows and trajectories in Lean~4 and uses verification feedback to refine them~\cite{wang2026lean4agent}. Both verify an \emph{interaction protocol or workflow of a particular task}, synthesized or checked per task. We verify a \emph{fixed, task-independent} coordination kernel once and govern the generated \emph{content}. The approaches compose: a verified per-task protocol could schedule the requests that our controller admits.

\emph{Validating individual LLM outputs.} Rigorous test-based evaluation reveals many incorrect programs that pass weaker tests~\cite{evalplus}, and model outputs vary across identical requests, even at temperature zero~\cite{ouyang2025nondeterminism}. RQ6 reproduces the first effect inside a coordinated workflow and attributes it to validator strength rather than to coordination, and RQ5 and RQ6 measure what the second means for early cutoff. Constrained decoding enforces syntactic and semantic constraints during generation~\cite{synchromesh,outlines}, and guardrail toolkits filter inputs and outputs~\cite{nemoguardrails}. All of these act on a single response. Our validators play the same role, but they are evaluated against the \emph{recorded} observation and combined with currentness, authority, structural and global checks. A response that passes every per-response check can still be stale, unauthorized or jointly inconsistent. Semantic caches reuse an earlier response for a similar request~\cite{bang2023gptcache}, and vCache bounds the resulting error rate~\cite{schroeder2026vcache}. Our early cutoff is the exact end of this spectrum: it reuses downstream work only when the regenerated value is equal, so it never admits a false hit.

\emph{Specifications for LLM systems.} Stoica et al.\ argue that specifications are the missing link for engineering LLM systems~\cite{stoica2024specifications}, and LLMs are increasingly used to produce them: SpecGen generates verifiable program specifications~\cite{ma2025specgen}, nl2postcond translates intent into postconditions~\cite{endres2024nl2postcond}, and recent evaluations find that LLM-written TLA$^+$ is often syntactically valid but rarely semantically correct~\cite{bisharat2026tlallm}. In our approach, users do not write TLA$^+$. They write jobs, read sets and validators, which are executable checks. The TLA$^+$ model concerns the kernel and is written once. Trace validation against TLA$^+$~\cite{cirstea2024trace,howard2025smartcasual} is the established path to keep a kernel implementation aligned with its model.

% ---------------------------------------------------------------------
\subsection{Foundations in Software Engineering, Databases and Programming Languages}
\emph{Consistency management and model transformation.} Instant consistency checking~\cite{egyed2006instant} and inconsistency tracking~\cite{egyed2011tracking} re-evaluate consistency rules on the model elements each rule accessed. Our read sets play the role of these scopes, extended with selected context and contracts. ATL~\cite{atl} provides typed transformation rules. Bidirectional transformations~\cite{foster2007lenses} and networks of models~\cite{stevens,klare} study consistency restoration across views and show that pairwise consistency does not imply global consistency, which is why we need a global predicate in addition to per-job validators. Stevens connects software build with megamodel consistency~\cite{megamodelbuild}. Our setting uses directed, single-producer stages with nondeterministic primitive proposals, and adds rejected values, obsolete replies and contract changes.

\emph{Incremental computation and build systems.} Self-adjusting computation~\cite{acar}, Adapton~\cite{hammer2014adapton}, incremental view maintenance~\cite{gupta1993dred}, build-system theory~\cite{mokhov}, Shake~\cite{mitchell2012shake} and Pluto~\cite{erdweg2015pluto} provide early cutoff and minimal recomputation for \emph{deterministic} tasks. Our keys are rebuild keys and our retained evidence is a one-entry cache. Determinism is what lets a build system treat a recorded trace as a proof that a rebuilt value is the right one. Without it, two valid repairs may differ, so we preserve admissibility relative to observations instead of equality with a fresh run, and failed-attempt memory adds an explicit unsuccessful outcome that build systems do not need.

\emph{Transactions and dynamic change.} OCC~\cite{occ} supplies the commit discipline, linearizability~\cite{herlihy1990linearizability} its correctness criterion, and the isolation anomalies of Berenson et al.~\cite{berenson1995critique} the vocabulary for what a version check does and does not exclude. Sagas~\cite{garcia1987sagas} replace atomicity by compensation, and CRDTs~\cite{shapiro2011crdt} avoid coordination for commutative updates. Neither applies directly: generated artifacts are not commutative, and our barrier makes compensation unnecessary. Kramer and Magee's quiescence~\cite{kramer1990evolving} and Vandewoude et al.'s tranquility~\cite{vandewoude2007tranquility} define safe dynamic change for components; Theorem~\ref{thm:reconfig} plays this role for agents, rules and validators.

\emph{Specification and refinement.} We use TLA$^+$~\cite{lamport1994tla,lamport2002specifying} with TLC~\cite{yu1999tlc} and TLAPS~\cite{cousineau2012tlaps}. Our refinement argument follows the forward-simulation method~\cite{lynch1995simulations,abadi1991refinement}, and IronFleet~\cite{ironfleet} shows how such refinements can be mechanized down to implementations. Our ablations follow the logic of mutation analysis~\cite{jia2011mutation}, applied to guards instead of statements.

% ---------------------------------------------------------------------
\subsection{What Is Borrowed and What Is New}
We borrow the medium-centric view of coordination and the artifact guard from tuple centres and A\&A, key-based binding from information protocols, the governor from electronic institutions, the commit discipline from OCC, and rebuild keys and early cutoff from build systems. Each of these mechanisms was designed for participants whose outputs are either trusted and deterministic (build systems, transactions) or untrusted but \emph{acting} rather than \emph{generating} (institutions, runtime enforcement). Untrusted, stochastic producers of interdependent artifacts break the assumptions of both groups. We see four contributions in combining these mechanisms for that setting:
(1) admissibility of \emph{content} relative to the observation it was computed from, with exact keys over everything that shapes the question (Theorem~\ref{thm:safety});
(2) early cutoff that remains sound when producers are nondeterministic, together with an exact cost bound (Theorems~\ref{thm:local} and~\ref{thm:cost});
(3) per-episode budgets and an upstream gate that turn ``eventually something happens'' into a terminal outcome and single-episode repair, including the configured-consumer structural rule that closes a liveness gap in our original design (Theorems~\ref{thm:single} and~\ref{thm:termination}, Lemma~\ref{lem:reactivation});
(4) a task-independent kernel whose guarantees are proved, model-checked with a refinement mapping, and partly machine-proved, so that users supply validators rather than specifications.
Section~\ref{sec:discussion} discusses, for each of these, what our evidence adds to what was already known.
